System identification is a fundamental step in the development of model-based control systems. Its purpose is to obtain mathematical models of dynamic systems from input-output data, allowing the behavior of a process to be represented without relying exclusively on first-principles equations.

In process control applications, these models are especially relevant because they provide the basis for simulation, controller design, performance analysis, and advanced process control (APC) strategies. A typical identification workflow involves selecting an appropriate model structure, estimating its parameters from data, validating its predictive performance, and assessing whether the model is adequate for the intended operating region. As stated by Box, all models are wrong, but some are useful~\cite{box1976science}; therefore, the objective is not to obtain a perfect representation of the process, but a model that captures the relevant dynamics with sufficient accuracy for control-oriented purposes.

This report focuses on the system identification of a nonlinear Continuous Stirred-Tank Reactor (CSTR) using simulated industrial-style step test data. The nonlinear reactor model is treated as the reference process, while local linear models are identified from its dynamic response to changes in the cooling water flow rate. This approach resembles the practical workflow used before implementing APC or model predictive control (MPC), where bump tests are performed to obtain dynamic models around a nominal operating point.

The CSTR model used in this work is based on Example 9-4, "Startup of a CSTR", from Fogler's \textit{Elements of Chemical Reaction Engineering}~\cite{fogler}. The model is derived from mass and energy balances and includes the effects of reaction kinetics, heat generation, and heat transfer through a cooling system. Although the original example focuses on reactor startup, this study uses the model to analyze local input-output dynamics around steady-state operation and to evaluate the suitability of different identified models for control-oriented applications.